% Volume VII: Formal Synthesis
% Part XIII. A Mathematics of Disciplined Distrust
% Chapter 74: Update Budgets
% Spine: proof-spines/ch074-spine.md

\chapter{Update Budgets}
\label{ch:ch074}

This chapter's thesis is that every responsible revision of belief should spend from an update budget whose size depends on the quality and reach of the incoming evidence. Drawing on the preceding account of premature-collapse loss and evidentiary reach, it proposes an informal rule: the better the provenance, independence, replication, representativeness, and adversarial survival of a record, the more movement it may justify. The budget idea clarifies why due process is not an ornamental delay but a mechanism for rationing how violently institutions may redraw a person's standing.
\[
\Theta_{t+1}^{\mathrm{free}}
=
rg\min_{\Theta} L(E_t;\Theta).
\]
\ToyModel\ An unconstrained institution would let new evidence drive whatever revision most reduces its local loss, regardless of how narrow or unrepresentative that evidence is.
It also explains why weak evidence can trigger attention without licensing wholesale reinterpretation of a case, a citizen, or a threat environment.
\[
D(\Theta_{t+1},\Theta_t)\le B(E_t).
\]
\ProposedTheorem\ Here \(D\) measures revision size and \(B(E_t)\) grows only with the tested quality and reach of the incoming record.
But the same rigor becomes an unmeetable burden when institutions define \(B(E_t)\) so parsimoniously for disfavored claimants that no attainable record can authorize meaningful correction. With revision now bounded, the next chapter can ask how privacy and publicity jointly shape the very evidence on which those budgets depend.
